\documentclass{article}
\usepackage{graphicx}
\usepackage[margin=1.5cm]{geometry}
\usepackage{amsmath}
\usepackage{array}
\usepackage{tikz}
\usetikzlibrary{arrows.meta,positioning,calc}

\begin{document}
\twocolumn

\title{Tuesday Review Document: Chapters 6--9}
\author{Prof. Jordan C. Hanson}
\date{}

\maketitle

\section*{Short Review Exercises}

\begin{enumerate}

\item \textbf{Adders and Arithmetic Logic}

A full-adder has inputs $A$, $B$, and $C_{\mathrm{in}}$, and outputs
$\Sigma$ and $C_{\mathrm{out}}$.

\begin{enumerate}
\item If $A=1$, $B=0$, and $C_{\mathrm{in}}=1$, determine
$\Sigma$ and $C_{\mathrm{out}}$.

\vspace{0.8cm}

\item Write the Boolean expression for the sum output $\Sigma$ in
terms of $A$, $B$, and $C_{\mathrm{in}}$.

\vspace{0.8cm}

\item Explain why several full-adders can be connected together to
add two multi-bit binary numbers.

\vspace{1.0cm}
\end{enumerate}

\item \textbf{Comparators, Decoders, and Encoders}

\begin{enumerate}
\item A 4-bit comparator receives
\[
A=1010, \qquad B=0111.
\]
Which comparator output is HIGH: $A>B$, $A=B$, or $A<B$?

\vspace{0.7cm}

\item A 4-to-16 decoder has input
\[
A_3A_2A_1A_0=1001.
\]
Which decimal output line should go HIGH?

\vspace{0.7cm}

\item A decimal-to-BCD encoder has the input line corresponding to
 decimal 6 HIGH. What 4-bit BCD code should appear at its output?

\vspace{0.7cm}
\end{enumerate}

\item \textbf{Multiplexers and Demultiplexers}

A 4-to-1 multiplexer has data inputs $D_0$ through $D_3$, select
inputs $S_1S_0$, and output $Y$.

\begin{enumerate}
\item Which data input is connected to $Y$ when
\[
S_1S_0=10?
\]

\vspace{0.6cm}

\item If
\[
D_0=0,\quad D_1=1,\quad D_2=1,\quad D_3=0,
\]
what is $Y$ when $S_1S_0=10$?

\vspace{0.6cm}

\item A 1-to-4 demultiplexer performs the reverse routing operation.
If $D_{\mathrm{in}}=1$ and $S_1S_0=01$, which output receives the
input bit?

\vspace{0.7cm}
\end{enumerate}

\item \textbf{Latches and Flip-Flops}

\begin{enumerate}
\item An SR latch initially has $Q=0$. If the SET input is asserted
momentarily and then released, what is the new value of $Q$?

\vspace{0.6cm}

\item A D flip-flop has $D=1$ immediately before its active clock
edge. What value is stored at $Q$ immediately after the edge?

\vspace{0.6cm}

\item A JK flip-flop has $J=1$ and $K=1$. If $Q=1$ before the active
clock edge, what is $Q$ afterward?

\vspace{0.7cm}
\end{enumerate}

\item \textbf{Shift Registers}

Consider a 4-bit serial-in/serial-out shift register whose flip-flops
are initially reset to 0.

\begin{enumerate}
\item A single 1 is shifted into the register, followed only by 0s.
How many active clock edges are required before the 1 reaches the
serial output?

\vspace{0.7cm}

\item If the clock frequency is $5\,\mathrm{MHz}$, how much time
passes before that bit reaches the output?

\vspace{0.8cm}

\item Briefly state the difference between \textit{serial loading}
and \textit{parallel loading} of a register.

\vspace{1.0cm}
\end{enumerate}

\item \textbf{Counters}

\begin{enumerate}
\item A 4-bit binary counter starts at $0000$. Write the next four
states after four active clock edges.

\vspace{0.8cm}

\item A 4-bit asynchronous counter is driven by a $100\,\mathrm{kHz}$
clock. How many clock cycles are required to count from $0000$ to
$1111$? How much time does this require?

\vspace{1.0cm}

\item What is the main difference between a synchronous counter and
an asynchronous (ripple) counter?

\vspace{1.0cm}
\end{enumerate}

\section*{Design Problem}

\item \textbf{Build a 4-bit Accumulator}

An \textbf{accumulator} is a digital system that repeatedly adds a new
binary value to a previously stored result. You will design a 4-bit
accumulator using components studied in class.

\begin{enumerate}

\item Begin with a \textbf{4-bit register} whose output is
\[
Q = Q_3Q_2Q_1Q_0.
\]
Draw the register and label the four data inputs, the four outputs, the
clock input \texttt{CLK}, and a reset input \texttt{RESET}. What value
should be stored in the register immediately after \texttt{RESET} is
activated?

\vspace{3.5cm}

\item Add a \textbf{4-bit parallel adder} to the system. One 4-bit
input to the adder is an external binary number
\[
X = X_3X_2X_1X_0.
\]

Connect the output of the register to the second 4-bit input of the
adder. Draw these connections and label the four sum outputs
\[
S = S_3S_2S_1S_0.
\]

Write an arithmetic expression relating $S$, $X$, and the current
register value $Q$.

\vspace{3.0cm}

\item Connect the four sum outputs of the adder to the four data
inputs of the register. Explain what happens on the next active clock
edge.

Write an expression for the next register value:
\[
Q_{\mathrm{next}} =
\]

\vspace{2.5cm}

\item Complete the feedback path so that the system can repeatedly
accumulate values. Explain why feeding the register output back into
the adder causes the system to remember the sum of all previous inputs.

\vspace{3.0cm}

\item Assume the accumulator is initially reset so that
\[
Q = 0000.
\]

The following values are presented to the input $X$, one value per
clock cycle:
\[
0011,\qquad 0010,\qquad 0100.
\]

Determine the value stored in the register after each clock edge.

\begin{center}
\renewcommand{\arraystretch}{1.25}
\begin{tabular}{c|c|c}
\hline
\textbf{Clock} & $\boldsymbol{X}$ & $\boldsymbol{Q}$ \textbf{ after edge} \\
\hline
0 & RESET & 0000 \\
1 & 0011 & \\
2 & 0010 & \\
3 & 0100 & \\
\hline
\end{tabular}
\end{center}

\item Suppose one additional input is
\[
X = 1000.
\]

Determine the new 4-bit register value. What happens if the numerical
result is larger than the largest value that can be stored in four
bits?

\vspace{2.5cm}

\end{enumerate}
\end{enumerate}

\end{document}
